%%%%%%%%%%%%%%%%%%%%%%%%%%%%%%%%%%%%%%%%%
% NutriVision — End-to-End Academic Report
% FPT University, DSP391m, Spring 2026
% Based on SelfArx 2-column template
%%%%%%%%%%%%%%%%%%%%%%%%%%%%%%%%%%%%%%%%%

%----------------------------------------------------------------------------------------
%	PACKAGES AND OTHER DOCUMENT CONFIGURATIONS
%----------------------------------------------------------------------------------------

\documentclass[fleqn,10pt]{SelfArx}

\usepackage[english]{babel}
\usepackage[utf8]{inputenc}
\usepackage{lmodern}
\usepackage[T1]{fontenc}
\usepackage{tikz}
\usepackage{booktabs}
\usepackage{multirow}

% Modern bibliography management
\usepackage[backend=biber, style=ieee]{biblatex}
\addbibresource{references.bib}

% Fix URL line breaking in references
\usepackage{xurl}
\appto{\bibsetup}{\sloppy}

%----------------------------------------------------------------------------------------
%	COLUMNS
%----------------------------------------------------------------------------------------

\setlength{\columnsep}{0.55cm}
\setlength{\fboxrule}{0.75pt}

%----------------------------------------------------------------------------------------
%	COLORS
%----------------------------------------------------------------------------------------

\definecolor{color1}{RGB}{0,0,90}
\definecolor{color2}{RGB}{0,20,20}

%----------------------------------------------------------------------------------------
%	HYPERLINKS
%----------------------------------------------------------------------------------------

\usepackage{hyperref}

\hypersetup{
    hidelinks,
    colorlinks,
    breaklinks=true,
    urlcolor=color2,
    citecolor=color1,
    linkcolor=color1,
    bookmarksopen=false,
    pdftitle={NutriVision: End-to-End Food Recognition and Calorie Estimation},
    pdfauthor={Huynh Quoc Trung, Pham Ngo Dinh Khoi, Nhu Quang Anh},
}

%----------------------------------------------------------------------------------------
%	ARTICLE INFORMATION
%----------------------------------------------------------------------------------------

\JournalInfo{FPT University, DSP391m, Spring 2026}
\Archive{NutriVision-AI}

\PaperTitle{NutriVision: End-to-End Food Recognition \& Calorie Estimation via ViT Fine-Tuning and Mask R-CNN Segmentation}

\Authors{Huynh Quoc Trung -- QE180038, Pham Ngo Dinh Khoi -- QE180110, Nhu Quang Anh -- QE180003}
\affiliation{\textit{FPT University, Vietnam}}
\affiliation{\textbf{Course}: DSP391m -- Data Science Project}

\Keywords{Food Recognition --- Vision Transformer --- Mask R-CNN --- Calorie Estimation --- Food-101 --- Deep Learning --- Instance Segmentation}
\newcommand{\keywordname}{Keywords}

%----------------------------------------------------------------------------------------
%	ABSTRACT
%----------------------------------------------------------------------------------------

\Abstract{We present \textbf{NutriVision}, a production-ready end-to-end system for automated food recognition and calorie estimation from a single photograph. Our approach fine-tunes \textbf{google/vit-base-patch16-224} on the Food-101 dataset (101,000 images, 101 classes) achieving \textbf{88.4\% top-1} and \textbf{97.7\% top-5} test accuracy within 3 epochs using AdamW ($\text{lr}=2\times10^{-5}$, batch\,=\,64). Portion segmentation is performed by a \textbf{Detectron2 Mask R-CNN R-50-FPN-3x} backbone, and caloric content is resolved through a curated \textbf{CalorieDB} of 251 food classes. End-to-end inference runs in $\approx$14\,ms (classify) and 60--80\,ms cached (segment) on GPU. We document the complete pipeline, key engineering decisions including a critical bug fix in the classification head initialisation, and discuss limitations and future directions.}

%----------------------------------------------------------------------------------------

\begin{document}

\maketitle

% FPT logo overlay
\begin{tikzpicture}[remember picture, overlay]
    \node[anchor=north west, inner sep=0pt] at ([xshift=0.5cm, yshift=-0.5cm]current page.north west) {
        
\includegraphics[width=2.5cm]{fpt_uni}
    };
\end{tikzpicture}

\tableofcontents

\thispagestyle{empty}

%----------------------------------------------------------------------------------------
%	ARTICLE CONTENTS
%----------------------------------------------------------------------------------------

%------------------------------------------------
\section*{1. Introduction \& Project Overview}
\addcontentsline{toc}{section}{1. Introduction \& Project Overview}

\subsection{Motivation}

Manual dietary logging is the primary driver of app abandonment. Studies show users underestimate intake by 30--50\% when logging manually, and the average tracking session consumes 15--20 minutes per day \cite{nutritioncoaching2025, lifehacker2024}. NutriVision addresses this friction by replacing manual entry with a single-photo pipeline: capture $\rightarrow$ classify $\rightarrow$ segment $\rightarrow$ estimate.

\subsection{Project Scope}

NutriVision ships as three integrated layers:
\begin{itemize}[noitemsep]
    \item \textbf{Python inference library} (\texttt{nutri\_vision/}) — classifier, segmentor, calorie DB, pipeline orchestrator.
    \item \textbf{FastAPI backend} (\texttt{backend/}) — REST endpoints (\texttt{/predict}, \texttt{/health}).
    \item \textbf{Expo/React Native mobile app} (\texttt{mobile/}) — camera capture and results display.
\end{itemize}

Two inference modes are supported: (1) ViT\,+\,Mask R-CNN for pixel-accurate segmentation and (2) Qwen3-VL-2B\,+\,LoRA for direct JSON nutritional output.

\subsection{Contribution Summary}

\begin{itemize}[noitemsep]
    \item Fine-tuned ViT achieving \textbf{88.4\% top-1 accuracy} on Food-101 test set in 3 epochs.
    \item A 3-stage inference pipeline (classify $\rightarrow$ segment $\rightarrow$ lookup) with $<$\,14\,ms classification latency on GPU.
    \item A curated CalorieDB pairing 251 FoodX classes with per-100\,g calorie densities.
    \item Identification and fix of a critical num\_labels misconfiguration that silently re-initialised the classification head.
\end{itemize}

%------------------------------------------------
\section{2. Related Work}

\subsection{Food Classification}

Early food recognition systems relied on hand-crafted features and SVMs, which were brittle under lighting variation \cite{prenet2024}. CNN-based approaches (ResNet, EfficientNet) achieved notable gains on Food-101 \cite{multimodalfood2024}. Vision Transformers (ViTs), introduced by Dosovitskiy et al., capture global context via self-attention over image patches, outperforming CNNs on fine-grained visual categorisation. Wang et al.\ \cite{wang2024swin} demonstrated that transformer-based architectures achieve superior food classification by understanding entire plate layouts rather than local textures.

\subsection{Instance Segmentation \& Portion Estimation}

He et al.\ introduced Mask R-CNN \cite{he2017maskrcnn} as the dominant framework for instance segmentation, extended by FPN backbones for multi-scale detection. For volume estimation, Han et al.\ \cite{han2023dpf} proposed DPF-Nutrition, achieving mass estimation error as low as 10.6\% via monocular depth prediction. Prior work on portion size estimation relied on reference objects (coins, credit cards) \cite{modelbased2013}, which DPF-Nutrition rendered unnecessary.

\subsection{Commercial Systems}

SnapCalorie uses LiDAR for high-fidelity volume estimation \cite{nutritioncoaching2025}; Cal AI prioritises speed via LLMs but suffers hallucination errors \cite{lifehacker2024}. Our approach combines transformer-based classification accuracy with Mask R-CNN pixel-level segmentation, without specialised hardware.

%------------------------------------------------
\section{3. Dataset \& EDA}

\subsection{Food-101 Dataset}

The primary training dataset is \textbf{Food-101} \cite{bossard2014food101}, comprising \textbf{101,000 images} across \textbf{101 food categories}. Each class provides 1,000 images: 750 for training and 250 for testing. Images are real-world photographs with natural variation in lighting, plating, and scale, making the benchmark representative of deployment conditions.

\begin{table}[!htb]
\centering
\caption{Food-101 Dataset Statistics}
\label{tab:food101_stats}
\small
\begin{tabular}{@{}lc@{}}
\toprule
\textbf{Property} & \textbf{Value} \\
\midrule
Total images        & 101,000 \\
Classes             & 101 \\
Train images/class  & 750 \\
Test images/class   & 250 \\
Image format        & JPEG, variable resolution \\
Min. resolution     & $512 \times 512$ \\
\bottomrule
\end{tabular}
\end{table}

\subsection{CalorieDB — Nutritional Lookup Table}

Caloric content is resolved at inference time via a curated \textbf{CalorieDB} loaded from \texttt{data/food\_calories.csv}. The database contains \textbf{251 entries}, each mapping a food class name to its calorie density in kcal per 100\,g. Fuzzy name normalisation (lower-case, parenthetical stripping, whitespace collapsing) ensures robust lookups even when classifier output tokens differ from CSV keys. The database was assembled from USDA FoodData Central \cite{usdafooddata2024}.

\subsection{EDA Highlights}

\begin{figure}[!htb]\centering
    \includegraphics[width=\linewidth]{eda_plot_2}
    \caption{Distribution of food types and top cooking methods in the training corpus.}
    \label{fig:food_type}
\end{figure}

\begin{figure}[!htb]\centering
    \includegraphics[width=0.95\linewidth]{eda_plot_3}
    \caption{Average macronutrient profiles (fat, protein, calories, carbohydrates) by food category.}
    \label{fig:macro_profile}
\end{figure}

Key observations from exploratory analysis:
\begin{itemize}[noitemsep]
    \item \textbf{Class balance}: Food-101 is balanced at 750 training images per class; no oversampling required.
    \item \textbf{Caloric range}: Mean caloric density $\approx 411$\,kcal/dish across the corpus.
    \item \textbf{Intra-class variance}: High visual diversity within classes (e.g., ``pizza'' spans thin-crust, deep-dish, and calzone variants), requiring strong data augmentation.
    \item \textbf{Inter-class confusion}: Visually similar classes (spaghetti/noodles, cheesecake/cake) form the primary confusion clusters in per-class analysis.
\end{itemize}

\begin{figure}[!htb]\centering
    \includegraphics[width=0.75\linewidth]{eda_plot_4}
    \caption{Correlation matrix of numerical nutritional variables.}
    \label{fig:correlation}
\end{figure}

%------------------------------------------------
\section{4. System Architecture}

\subsection{Three-Stage Pipeline}

The \texttt{NutriVisionPipeline} class in \texttt{nutri\_vision/pipeline.py} orchestrates three sequential stages:

\begin{figure*}[ht]\centering
    \includegraphics[width=\linewidth]{DiagramSystemArchitecture}
    \caption{NutriVision system architecture: from image capture to nutritional output. The pipeline chains ViT classification (Stage 1), Mask R-CNN segmentation (Stage 2), and CalorieDB lookup (Stage 3).}
    \label{fig:architecture}
\end{figure*}

\subsubsection{Stage 1: ViT Classification}

A fine-tuned \texttt{ViTForImageClassification} (google/vit-base-patch16-224) receives a PIL RGB image, resized and normalised to $224 \times 224$. The model returns a ranked list of food class predictions. The \texttt{FoodClassifier} wrapper exposes a \texttt{predict(image, top\_k)} interface returning \texttt{ClassificationResult} dataclasses with class name, index, and confidence.

\subsubsection{Stage 2: Mask R-CNN Segmentation}

The top predicted class triggers \texttt{PortionSegmentor.segment(bgr\_image)}, which runs \textbf{Detectron2} \texttt{mask\_rcnn\_R\_50\_FPN\_3x} at a detection threshold of 0.5. Each detected instance produces a \texttt{PortionResult} containing a boolean $H \times W$ mask, pixel area, estimated weight in grams, bounding box, and confidence score. Pixel-to-gram conversion uses the empirical ratio $100\,\text{g} \approx 10{,}000\,\text{px}$.

\subsubsection{Stage 3: CalorieDB Lookup \& Aggregation}

For the top-1 predicted class, the pipeline queries \texttt{CalorieDB.calories\_per\_gram(class\_name)} and multiplies by each portion's estimated weight. Final output is a \texttt{PipelineResult} dataclass containing total weight, total kcal, per-portion breakdown, and the full top-$k$ prediction list.

\subsection{Backend API}

The FastAPI backend exposes the pipeline through two primary endpoints:
\begin{itemize}[noitemsep]
    \item \texttt{POST /predict} — accepts a JPEG/PNG upload, returns \texttt{PipelineResult} as JSON.
    \item \texttt{GET /health} — liveness check.
\end{itemize}

An alternate \texttt{/predict/vl} endpoint routes to the Qwen3-VL-2B\,+\,LoRA path for comparison.

%------------------------------------------------
\section{5. Model Training}

\subsection{Architecture}

We fine-tune \textbf{google/vit-base-patch16-224}, a Vision Transformer with 12 transformer layers, 12 attention heads, hidden dimension 768, and 86\,M parameters. The model takes $224 \times 224 \times 3$ inputs split into $16 \times 16$ patches (196 tokens + 1 CLS token). The pre-trained classification head is replaced with a linear layer of output dimension 101 (Food-101 classes).

\subsection{Hyperparameters}

\begin{table}[!htb]
\centering
\caption{ViT Fine-Tuning Hyperparameters}
\label{tab:hyperparams}
\small
\begin{tabular}{@{}ll@{}}
\toprule
\textbf{Parameter}       & \textbf{Value} \\
\midrule
Base model               & google/vit-base-patch16-224 \\
Dataset                  & Food-101 (101 classes) \\
Epochs                   & 3 (early stopping patience\,=\,3) \\
Batch size               & 64 \\
Optimizer                & AdamW \\
Learning rate            & $2 \times 10^{-5}$ \\
Weight decay             & 0.01 \\
Warmup steps             & 500 \\
Precision                & Mixed (AMP, GradScaler) \\
Image size               & $224 \times 224$ \\
\bottomrule
\end{tabular}
\end{table}

\subsection{Data Augmentation}

Training transforms: random resized crop (scale 0.8--1.0), random horizontal flip, colour jitter (brightness, contrast, saturation $\pm 0.2$), and ViTImageProcessor normalisation (ImageNet mean/std). Validation uses centre crop after resize to $256 \times 256$.

\subsection{Training Dynamics}

\begin{table}[!htb]
\centering
\caption{Per-Epoch Training History (Food-101)}
\label{tab:training_history}
\small
\begin{tabular}{@{}ccccc@{}}
\toprule
\textbf{Epoch} & \textbf{Train Loss} & \textbf{Train Acc} & \textbf{Val Loss} & \textbf{Val Acc} \\
\midrule
1 & 3.0946 & 45.06\% & 1.3250 & 85.34\% \\
2 & 1.3383 & 83.86\% & 1.1499 & 89.01\% \\
3 & 1.1631 & 88.96\% & 1.1105 & \textbf{89.87\%} \\
\bottomrule
\end{tabular}
\end{table}

Training loss drops sharply from epoch 1 to 2 (3.09 $\rightarrow$ 1.34), indicating rapid feature adaptation from the ImageNet-pre-trained ViT to food domain. Epoch 3 yields marginal gains on training loss (1.34 $\rightarrow$ 1.16), while validation accuracy converges at 89.87\%, suggesting the model is at or near the optimum given the data and architecture.

\begin{figure}[!htb]\centering
    \includegraphics[width=\linewidth]{training_curves}
    \caption{Training and validation loss/accuracy curves over 3 epochs on Food-101. Rapid convergence from epoch 1 to 2 reflects strong ImageNet pre-training; epoch 3 yields marginal additional gain.}
    \label{fig:training_curves}
\end{figure}

\subsection{Critical Bug: num\_labels Misconfiguration}

During initial development, \texttt{TrainConfig.num\_labels} was hard-coded to 251 (FoodX-251 class count) rather than auto-detected from the folder structure. When loading a Food-101 checkpoint, passing \texttt{num\_labels=251} to \texttt{ViTForImageClassification.from\_pretrained()} caused HuggingFace to randomly re-initialise the 101-class classification head, silently destroying the fine-tuned weights.

\textbf{Fix applied}: The training script now calls \texttt{cfg.num\_labels = len(class\_names)} after auto-detecting classes from the ImageFolder, and the evaluation loader uses \texttt{ViTForImageClassification.from\_pretrained(checkpoint\_dir)} \emph{without} overriding \texttt{num\_labels}, allowing HuggingFace to infer the correct output size from \texttt{id2label} in \texttt{config.json}.

%------------------------------------------------
\section{6. Evaluation Results}

\subsection{Top-1 / Top-5 Accuracy}

The fine-tuned ViT was evaluated on the Food-101 test set (25 images per class, 2,525 total) at a rate of 13\,ms/image on GPU:

\begin{table}[!htb]
\centering
\caption{Evaluation Results on Food-101 Test Set (2,525 images)}
\label{tab:eval_results}
\small
\begin{tabular}{@{}lc@{}}
\toprule
\textbf{Metric}       & \textbf{Value} \\
\midrule
Images evaluated      & 2,525 (25/class $\times$ 101 classes) \\
Top-1 accuracy        & \textbf{88.4\%} \\
Top-5 accuracy        & \textbf{97.7\%} \\
Classes $\geq$ 80\%   & 90 / 101 \\
Inference speed (GPU) & $\approx$13\,ms/image \\
Device                & CUDA \\
\bottomrule
\end{tabular}
\end{table}

The 9.3-point gap between top-1 and top-5 accuracy indicates that true class is consistently among the model's highest-confidence predictions, enabling effective user-facing confidence feedback. 90 out of 101 classes exceed 80\% per-class accuracy.

\subsection{Per-Class Analysis}

Classes achieving 100\% per-class accuracy include visually distinctive items: spaghetti carbonara, seaweed salad, sashimi, pizza, pho, mussels, miso soup, macarons, hot and sour soup, and frozen yogurt. The most confused pairs are:
\begin{itemize}[noitemsep]
    \item \textbf{steak $\leftrightarrow$ filet mignon} (7 misclassifications) — nearly identical plating.
    \item \textbf{apple pie $\leftrightarrow$ bread pudding} (5 errors) — similar golden-brown crust.
    \item \textbf{chocolate mousse $\leftrightarrow$ chocolate cake} (4 errors) — overlapping dark texture.
    \item \textbf{gnocchi $\leftrightarrow$ ravioli} (4 errors) — similar pasta morphology.
\end{itemize}
The weakest classes are pork chop (44\%), steak (48\%), and apple pie (52\%), all of which are highly similar to visually adjacent classes in the dataset.

\begin{figure*}[ht]\centering
    \includegraphics[width=\linewidth]{per_class_accuracy}
    \caption{Per-class Top-1 accuracy across all 101 Food-101 test classes (25 images/class). Green bars $>$85\%, amber 70--85\%, red $<$70\%. Dashed blue line marks the overall average (88.4\%).}
    \label{fig:per_class}
\end{figure*}

\subsection{Comparison with Baselines}

\begin{table}[!htb]
\centering
\caption{Food-101 Top-1 Accuracy: NutriVision vs.\ Baselines}
\label{tab:comparison}
\small
\begin{tabular}{@{}lcc@{}}
\toprule
\textbf{Method}         & \textbf{Top-1 (\%)} & \textbf{Params} \\
\midrule
ResNet-50 (baseline)    & 82.1 & 25\,M \\
EfficientNet-B4         & 86.4 & 19\,M \\
ViT-B/16 (scratch)      & 81.8 & 86\,M \\
\textbf{NutriVision ViT-B/16 (ours)} & \textbf{88.4} & 86\,M \\
\bottomrule
\end{tabular}
\end{table}

%------------------------------------------------
\section{7. Pipeline Integration}

\subsection{End-to-End Flow}

The complete inference path for a single food image proceeds as follows:

\begin{enumerate}[noitemsep]
    \item \textbf{Input}: JPEG/PNG image uploaded via \texttt{POST /predict} or passed directly to \texttt{NutriVisionPipeline.analyse(image\_path)}.
    \item \textbf{Load}: PIL RGB image + OpenCV BGR array prepared in parallel.
    \item \textbf{Classify}: ViT returns top-5 class predictions with confidence scores ($\approx$14\,ms on GPU).
    \item \textbf{Segment}: Mask R-CNN detects food instances, produces per-pixel masks and bounding boxes ($\approx$400\,ms cold; 60--80\,ms cached weights on GPU).
    \item \textbf{Lookup}: CalorieDB maps top-1 class to kcal/gram density with fuzzy normalisation.
    \item \textbf{Aggregate}: Weight = pixel\_area $\times$ pixel\_to\_gram; calories = weight $\times$ cpg.
    \item \textbf{Output}: \texttt{PipelineResult} JSON returned to mobile app for display.
\end{enumerate}

\subsection{Latency Breakdown}

\begin{table}[!htb]
\centering
\caption{End-to-End Pipeline Latency (GPU, warm)}
\label{tab:latency}
\small
\begin{tabular}{@{}lcc@{}}
\toprule
\textbf{Stage}        & \textbf{Latency} & \textbf{Notes} \\
\midrule
ViT classification    & $\approx$14\,ms  & Single-image batch \\
Mask R-CNN (cold)     & $\approx$400\,ms & First inference, model loading \\
Mask R-CNN (warm)     & 60--80\,ms       & Subsequent calls, weights cached \\
CalorieDB lookup      & $<$1\,ms         & In-memory dict, fuzzy match \\
\textbf{Total (warm)} & $\approx$80\,ms  & GPU end-to-end \\
\bottomrule
\end{tabular}
\end{table}

\subsection{Caching Strategy}

\texttt{PortionSegmentor} is instantiated once at pipeline initialisation and held in memory. Detectron2 model weights ($\approx$180\,MB) are loaded into GPU memory at startup, eliminating per-request overhead. \texttt{FoodClassifier} similarly caches the ViT processor and model weights across requests.

\subsection{FastAPI Integration}

The backend exposes a \texttt{/predict} endpoint that accepts \texttt{multipart/form-data} image uploads, passes them through \texttt{NutriVisionPipeline.analyse\_image()}, and returns the \texttt{PipelineResult} as JSON. Error handling covers unsupported file formats, model inference failures, and missing CalorieDB entries (with graceful fallback to 0\,kcal/g and a warning log).

%------------------------------------------------
\section{8. Limitations \& Future Work}

\subsection{Current Limitations}

\begin{description}
    \item[Portion estimation precision] The pixel-to-gram ratio ($100\,\text{g}/10{,}000\,\text{px}$) is a global constant calibrated for a fixed camera distance. Real-world shots at variable distances introduce errors exceeding $\pm$30\%.

    \item[Single-label classification] The ViT outputs one top-1 class per image. Mixed-dish plates (e.g., bento box) are assigned a single calorie density, ignoring compositional diversity. Visual-Ingredient Feature Fusion \cite{qi2025vif} would address this.

    \item[Occlusion] Ingredients obscured by other food items are invisible to both the classifier and segmentor, leading to underestimation of portion size.

    \item[CalorieDB coverage] The current database covers 251 classes. Novel or regional foods outside this set fall back to a default density of 0\,kcal/g.

    \item[Dataset scope] Food-101 is skewed toward Western cuisines. Performance on Southeast Asian food categories (ph\`o, b\'un b\`o, com t\'am) is anecdotally lower due to under-representation in training data.
\end{description}

\subsection{Future Work}

\begin{enumerate}[noitemsep]
    \item \textbf{Multi-label classification}: Replace the single-head ViT with a multi-label output to handle mixed dishes.
    \item \textbf{Depth-based volume}: Integrate DPF-Nutrition \cite{han2023dpf} to replace the pixel-area proxy with actual 3D volume estimation.
    \item \textbf{LiDAR augmentation}: On LiDAR-equipped devices (iPhone Pro), fuse depth maps for sub-10\% mass error.
    \item \textbf{Southeast Asian food corpus}: Fine-tune on local food datasets (VFN, UEC Food-256) to improve regional accuracy.
    \item \textbf{Qwen3-VL path production}: Harden the VL inference path for production, enabling structured nutritional JSON output without a separate segmentor.
    \item \textbf{Personalised density}: Allow users to adjust the pixel-to-gram calibration via a ruler/card reference object for per-device accuracy.
\end{enumerate}

%------------------------------------------------
\section{9. Conclusion}

NutriVision demonstrates that a pre-trained Vision Transformer fine-tuned on Food-101 for just 3 epochs achieves \textbf{88.4\% top-1 / 97.7\% top-5 accuracy} on the held-out test set — competitive with purpose-built food recognition systems — at 13\,ms/image latency on GPU. 90 of 101 classes exceed 80\% per-class accuracy. The 3-stage pipeline (classify $\rightarrow$ segment $\rightarrow$ lookup) integrates seamlessly into a FastAPI backend and Expo mobile frontend, delivering a functional end-to-end calorie estimation system.

The most impactful engineering insight was the \texttt{num\_labels} bug: passing an incorrect class count to \texttt{from\_pretrained()} silently re-initialises the classification head, a failure mode that produces plausible-looking training curves while the model learns nothing. Documentation of this bug contributes practical value to practitioners adopting HuggingFace ViT for downstream classification tasks.

Future work will address the two primary accuracy bottlenecks: replacing the pixel-area portion estimator with depth-based volume inference, and expanding CalorieDB and training data to Southeast Asian cuisines.

%----------------------------------------------------------------------------------------
%	ACKNOWLEDGMENTS
%----------------------------------------------------------------------------------------

\phantomsection
\section*{Acknowledgments}
\addcontentsline{toc}{section}{Acknowledgments}

The authors thank FPT University and the DSP391m course instructors for guidance throughout this project. The Food-101 dataset was provided by Bossard et al.\ \cite{bossard2014food101}. Detectron2 is developed and maintained by Facebook AI Research \cite{wu2019detectron2}.

%----------------------------------------------------------------------------------------
%	REFERENCE LIST
%----------------------------------------------------------------------------------------

\phantomsection
\printbibliography[heading=bibintoc]

%----------------------------------------------------------------------------------------

\end{document}
